\subsection{Using wrappers vs.~scenario calls vs.~direct calls}
\label{sec:wrapper_vs_direct}

\begin{lstlisting}[language=Python, caption={\textbf{Vacuum and constant density, three ways.} The three-flavor probability, computed three ways. The wrappers take the oscillation parameters as keywords, here the NuFIT 6.1 values, and the density in g~cm$^{-3}$. The scenario function takes the density as a function of position. The hand-built form is the vacuum term divided by the energy plus the potential $V_{\rm CC}$ times the projector of \equ{h_matter}; {\tt VCC\_EARTH\_CRUST} is that potential at 3~g~cm$^{-3}$ and $Y_e = 1/2$. All three return the same probabilities; the times are measured under the protocol of Sec.~\ref{sec:timing_protocol}. See Sec.~\ref{sec:wrapper_vs_direct} for details.}, label={lst:constant}, float=tbp]
import numpy as np
import magnus.globaldefs as gd
import magnus.hamiltonians as hamiltonians
import magnus.matter as matter
import magnus.oscprob as oscprob

osc = gd.load_nufit_params('NuFIT 6.1')
E, L = 1.0*gd.UNIT_GEV, 1000.0*gd.UNIT_KM
Es = np.logspace(-1, 1, 200)*gd.UNIT_GEV

# --- 1. With the named wrappers
P_vac = oscprob.osc_prob_3nu_vacuum(E, L, **osc)
P_mat = oscprob.osc_prob_3nu_matter_constant_density(
    E, L, rho=3.0, density_matter_is_in_g_per_cm3=True, **osc)
# Pme = 0.00377306 in vacuum, 0.01347547 in
# matter, in 0.090 ms

P_scan = oscprob.osc_prob_3nu_matter_constant_density(
    Es, L, rho=3.0, density_matter_is_in_g_per_cm3=True, **osc)
# P_scan.shape = (200, 3, 3), in 0.123 ms

# --- 2. With a scenario function
rho_func = lambda l: 3.0*gd.UNIT_G_PER_CM3
P_mat = oscprob.osc_prob_matter_std_potential(
    3, rho_func, E, L, osc_params=osc)
# The same number, to 4e-16, in 31.9 ms

# --- 3. With a Hamiltonian built by hand
H_vac = hamiltonians.hamiltonian_3nu_vacuum_energy_independent(
    s12=osc['s12'], s23=osc['s23'], s13=osc['s13'], dCP=osc['dCP'],
    D21=osc['D21'], D31=osc['D31'])
proj = matter.matter_potential_projector(3)

P_vac = oscprob.osc_prob(H_vac/E, 0.0, L)
P_mat = oscprob.osc_prob(
    H_vac/E + gd.VCC_EARTH_CRUST*proj, 0.0, L)
# The same two numbers, in 0.013 ms

P_scan = oscprob.osc_prob_energy_baseline(
    lambda e: H_vac/e + gd.VCC_EARTH_CRUST*proj, Es, L,
    H_func_is_function_only_of_energy=True)
# The same 200 points, in 3.6 ms
\end{lstlisting}
%%%%%%%%%%%%%%%%%%%%%%%%%%%%%%%%%%%%%%%%%%%%%%%%%%%%%

Every wrapper function (\eg, {\tt osc\_prob\_3nu\_vacuum}, {\tt osc\_prob\_3nu\_matter\_constant\_density}) builds a Hamiltonian internally and hands it, through a scenario function, to the engines of Sec.~\ref{sec:engines}, the last of which is {\tt osc\_prob}, the primitive probability function of \magnus. The wrappers are a convenience; a user can do the same by hand, or enter at the scenario layer between the two. Wrappers also contain the most safeguards and validation particular to the cases they are built for; scenario functions contain fewer; and the direct call to {\tt osc\_prob}, the least.  Figure~\ref{fig:layers} shows the list of wrappers, scenario functions, and their relation to one another and {\tt osc\_prob}.

Listing~\ref{lst:constant} computes the same three-flavor probability all three ways, at one energy and over a scan of two hundred. The three agree: to round-off in matter, bit for bit in vacuum. 

%%%%%%%%%%%%%%%%%%%%%%%%%%%%%%%%%%%%%%%%%%%%%%%%%%%%%

\subsubsection{Through a wrapper}
\label{sec:wrappers}

%%%%%%%%%%%%%%%%%%%%%%%%%%%%%%%%%%%%%%%%%%%%%%%%%%%%%
{\def\arraystretch{1.12}
\begin{table*}[t!]
{\footnotesize
\begin{tabular}{@{}p{2.7cm}p{2.5cm}p{5.9cm}p{5.9cm}@{}}
 Environment & Physics & \multicolumn{2}{@{}l@{}}{Wrapper function} \\
 \hline
 Vacuum & Standard & {\tt osc\_prob\_2nu\_vacuum} & {\tt osc\_prob\_3nu\_vacuum} \\
  &  & {\tt osc\_prob\_4nu\_vacuum} & {\tt osc\_prob\_5nu\_vacuum} \\
  & Lorentz violation & {\tt osc\_prob\_2nu\_vacuum\_liv} & {\tt osc\_prob\_3nu\_vacuum\_liv} \\
  &  & {\tt osc\_prob\_4nu\_vacuum\_liv} & {\tt osc\_prob\_5nu\_vacuum\_liv} \\
 Constant density & Standard & {\tt osc\_prob\_2nu\_matter\_constant\_density} & {\tt osc\_prob\_3nu\_matter\_constant\_density} \\
  &  & {\tt osc\_prob\_4nu\_matter\_constant\_density} & {\tt osc\_prob\_5nu\_matter\_constant\_density} \\
  & Non-standard int. & {\tt osc\_prob\_2nu\_matter\_nsi\_constant\_density} & {\tt osc\_prob\_3nu\_matter\_nsi\_constant\_density} \\
  &  & {\tt osc\_prob\_4nu\_matter\_nsi\_constant\_density} & {\tt osc\_prob\_5nu\_matter\_nsi\_constant\_density} \\
  & Lorentz violation & {\tt osc\_prob\_2nu\_matter\_liv\_constant\_density} & {\tt osc\_prob\_3nu\_matter\_liv\_constant\_density} \\
  &  & {\tt osc\_prob\_4nu\_matter\_liv\_constant\_density} & {\tt osc\_prob\_5nu\_matter\_liv\_constant\_density} \\
 Exponential density & Standard & {\tt osc\_prob\_2nu\_matter\_exp\_density} & {\tt osc\_prob\_3nu\_matter\_exp\_density} \\
  &  & {\tt osc\_prob\_4nu\_matter\_exp\_density} & {\tt osc\_prob\_5nu\_matter\_exp\_density} \\
  & Non-standard int. & {\tt osc\_prob\_2nu\_matter\_nsi\_exp\_density} & {\tt osc\_prob\_3nu\_matter\_nsi\_exp\_density} \\
  &  & {\tt osc\_prob\_4nu\_matter\_nsi\_exp\_density} & {\tt osc\_prob\_5nu\_matter\_nsi\_exp\_density} \\
  & Lorentz violation & {\tt osc\_prob\_2nu\_matter\_liv\_exp\_density} & {\tt osc\_prob\_3nu\_matter\_liv\_exp\_density} \\
  &  & {\tt osc\_prob\_4nu\_matter\_liv\_exp\_density} & {\tt osc\_prob\_5nu\_matter\_liv\_exp\_density} \\
 Earth (PREM) & Standard & {\tt osc\_prob\_2nu\_earth} & {\tt osc\_prob\_3nu\_earth} \\
  &  & {\tt osc\_prob\_4nu\_earth} & {\tt osc\_prob\_5nu\_earth} \\
  & Non-standard int. & {\tt osc\_prob\_2nu\_earth\_nsi} & {\tt osc\_prob\_3nu\_earth\_nsi} \\
  &  & {\tt osc\_prob\_4nu\_earth\_nsi} & {\tt osc\_prob\_5nu\_earth\_nsi} \\
  & Lorentz violation & {\tt osc\_prob\_2nu\_earth\_liv} & {\tt osc\_prob\_3nu\_earth\_liv} \\
  &  & {\tt osc\_prob\_4nu\_earth\_liv} & {\tt osc\_prob\_5nu\_earth\_liv} \\
 Sun & Standard & {\tt osc\_prob\_2nu\_sun} & {\tt osc\_prob\_3nu\_sun} \\
  &  & {\tt osc\_prob\_4nu\_sun} & {\tt osc\_prob\_5nu\_sun} \\
  & Non-standard int. & {\tt osc\_prob\_2nu\_sun\_nsi} & {\tt osc\_prob\_3nu\_sun\_nsi} \\
  &  & {\tt osc\_prob\_4nu\_sun\_nsi} & {\tt osc\_prob\_5nu\_sun\_nsi} \\
  & Lorentz violation & {\tt osc\_prob\_2nu\_sun\_liv} & {\tt osc\_prob\_3nu\_sun\_liv} \\
  &  & {\tt osc\_prob\_4nu\_sun\_liv} & {\tt osc\_prob\_5nu\_sun\_liv} \\
\end{tabular}}
\caption{\label{tab:wrappers}\textbf{The probability wrapper functions.} The 56 wrapper functions of \magnus, by environment and by the physics in the Hamiltonian. Each of the fourteen combinations exists at two, three, four and five flavors. There is no wrapper for non-standard interactions in vacuum, since those are interactions with matter. Every function here builds a Hamiltonian and hands it to the engines of Sec.~\ref{sec:engines}. (Three others, not shown here, do not: {\tt osc\_prob\_2nu\_vacuum\_std}, {\tt osc\_prob\_3nu\_vacuum\_std} and {\tt osc\_prob\_2nu\_matter\_std} evaluate the standard analytic expressions for the probability instead. See Sec.~\ref{sec:wrappers} for details.}
\end{table*}}
%%%%%%%%%%%%%%%%%%%%%%%%%%%%%%%%%%%%%%%%%%%%%%%%%%%%%

\tabl{wrappers} lists the 56 wrappers shipped with \magnus. A wrapper is the shortest call whenever the scenario is one \magnus\ already knows. It takes the oscillation parameters as keywords, fills in any left unset from a named set, converts the density from g~cm$^{-3}$, and assembles the Hamiltonian. That last step includes the vacuum term: {\tt osc\_prob\_3nu\_matter\_constant\_density} takes the same six oscillation parameters as {\tt osc\_prob\_3nu\_vacuum} and internally adds the potential to what it builds from them, which is why {\tt H\_vac} appears only in the third block of Listing~\ref{lst:constant}.

Every wrapper takes its arguments in the same four groups: the energy, the geometry of the environment, the oscillation parameters, and whatever the scenario adds. For instance, for non-standard interactions (NSI) at three flavors in constant density,
\begin{pysnip}
P = oscprob.\
 osc_prob_3nu_matter_nsi_constant_density(
    E, L,            # 1. energy, baseline
    rho=3.0,         # 2. density
    density_matter_is_in_g_per_cm3=True,
    **osc,           # 3. mixing
    eps_em=0.05)     # 4. new physics
# Pme = 0.016972.  Dropping eps_em gives
# 0.013475, the standard wrapper's answer
\end{pysnip}
The fourth group (``new physics'') exists only in the wrappers whose name carries {\tt \_nsi} or {\tt \_liv}; \tabl{parameters} lists its parameters. Every one of them defaults to zero, so one of those wrappers called without them returns the standard result, bit for bit identical to what the standard wrapper of the same environment returns. That makes it the control for any new-physics scan. 

Passing {\tt rho} as a scalar makes the Hamiltonian independent of position, so the request reaches the constant-Hamiltonian engine of Sec.~\ref{sec:engines} and the whole scan becomes one stack of matrix exponentials. Two hundred energies cost 0.123~ms vs.~0.090~ms for one, so the scan is nearly free.

%%%%%%%%%%%%%%%%%%%%%%%%%%%%%%%%%%%%%%%%%%%%%